% ======================================================================
\section{Summing over paths}\label{sec:counting}
% ======================================================================

By \cref{lem:selection,prop:cost}, the susceptibility is bounded by sums, over paths $P$ from a fixed root, of $u^{|P|}e^{\theta c(P)}$ times shape factors. This section bounds such sums uniformly in $L$. The two ingredients are an exact moment of the chord number (\cref{lem:rosenbluth}) and transfer bounds on a finite automaton that remembers the last $16$ steps (\cref{lem:transfer,cert:transfer}). H\"older's inequality combines them (\cref{lem:holder}). Fix an even $L\ge8$ and a root $x\in\T_L$. All paths below start at $x$.

\subsection{Contacts}

For a path $P=(v_0,\dots,v_n)$ and $1\le i\le n$ let
\[
 a_i(P)=\#\{j\le i-2:\ v_j\sim v_i\}
\]
be the number of \emph{contacts} of the tip $v_i$ with earlier vertices other than its predecessor. For $0\le i\le n$ let $B_i(P)$ be the set of directions $d\in\{\pm e_1,\pm e_2,\pm e_3\}$ such that $v_i+d\notin\{v_0,\dots,v_i\}$.

\begin{lemma}\label{lem:contacts}
$c(P)=\sum_{i=1}^na_i(P)$, and $|B_i(P)|=5-a_i(P)$ for $1\le i\le n$. Moreover $|B_0(P)|=6$, and $v_{i+1}-v_i\in B_i(P)$ for $i<n$.
\end{lemma}

\begin{proof}
A chord $\{v_j,v_i\}$ with $j\le i-2$ is counted exactly once, in $a_i$. The six neighbours of $v_i$ are distinct, since $L\ge3$. One of them is $v_{i-1}$, $a_i$ of them are earlier vertices of $P$, and the others are unvisited at time $i$.
\end{proof}

\subsection{An exact chord moment}

\begin{lemma}[Chord moment]\label{lem:rosenbluth}
For every $n\ge1$,
\[
 \sum_{|P|=n}\Bigl(\frac54\Bigr)^{c(P)}\ \le\ 6\Bigl(\frac54\Bigr)^5 5^{n-1}.
\]
\end{lemma}

\begin{proof}
Put $b=5/4$. The five integer inequalities $(5-a)b^a\le5$, $a=0,\dots,5$, are checked directly. The cases $a=0,1$ are equalities, and the values for $a=2,3,4,5$ are $4.6875$, $3.90625$, $2.44140625$ and $0$. For a path $p$ of length $\ell\ge1$ and $k\ge0$ let
\[
 F_k(p)=\sum_{q}\ \prod_{i=\ell}^{\ell+k-1}b^{a_i(q)},
\]
the sum running over all paths $q$ of length $\ell+k$ extending $p$. Then $F_0=1$, and by \cref{lem:contacts}
\[
 F_{k+1}(p)=b^{a_\ell(p)}\sum_{d\in B_\ell(p)}F_k(p+d)\le b^{a_\ell(p)}(5-a_\ell(p))\max_dF_k(p+d).
\]
By induction $F_k(p)\le5^k$. The root has six continuations, and the final tip contributes $b^{a_n}\le b^5$. Hence $\sum_{|P|=n}b^{c(P)}\le b^5\cdot6\cdot5^{n-1}$.
\end{proof}

The argument is the Rosenbluth growth construction~\cite{RR55} read as a counting inequality. The weight $\prod_i(5-a_i)^{-1}$ is the probability that a walk choosing uniformly among unvisited neighbours produces $P$. The lemma has been verified in Lean on $(\Z/L\Z)^3$ for every $L\ge3$ (Appendix~\ref{app:lean}).

\subsection{A finite-memory automaton}

Self-avoidance is a global constraint. We replace it by a local one: a step is forbidden only if it returns to one of the recently visited vertices that are still close to the tip. Paths on $\T_L$ are first lifted to $\Z^3$. The \emph{lift} of $P$ is the lattice walk $\omega=(x_0,\dots,x_n)$ with $x_0=0$ and the same step directions $d_i=v_i-v_{i-1}$. Distinct sites of $\T_L$ have distinct lifts, so $\omega$ is self-avoiding.

\begin{definition}[Suffix automaton]\label{def:automaton}
Fix the memory $M=16$.
\begin{enumerate}
\item A \emph{retained suffix} is a sequence $(y_0,\dots,y_m)$ of distinct sites of $\Z^3$ with $y_k\sim y_{k-1}$ for all $k$.
\item For a direction $d$ with $y_m+d\notin\{y_0,\dots,y_m\}$, the \emph{update} by $d$ appends $y_{m+1}=y_m+d$. It then deletes $y_0,\dots,y_{j-1}$, where $j$ is the smallest index with $(m+1-j)+\|y_{m+1}-y_j\|_1\le M$.
\item A \emph{state} is a retained suffix modulo translations and the $48$ signed permutations of the axes. The \emph{initial state} $s_0$ is the one-point suffix at the root.
\item For a state $s$, the \emph{available labels} $A(s)$ are the directions $d$ with $y_m+d$ outside the suffix, and $s_d$ is the state after the update by $d$. Labels are counted with multiplicity, also when different labels lead to the same state.
\end{enumerate}
The \emph{automaton} is the finite set of states reachable from $s_0$, together with the labelled transitions $s\to s_d$.
\end{definition}

A vertex $y_j$ is kept as long as a self-avoiding continuation could return to it within $M$ steps. Only two properties of the rule are used below. First, the retained suffix consists of actual vertices of the walk. Second, after at least two steps it contains the last three vertices, since $2+\|y_m-y_{m-2}\|_1\le4\le M$. Consequently the state determines the last two step directions, and with them the shape (\cref{def:shape}) of the triple completed by the next step.

\begin{lemma}[Domination by the automaton]\label{lem:automaton}
Let $P$ be a path on $\T_L$ and let $s_0,s_1,\dots,s_n$ be the states obtained from the initial state by feeding the step directions of $P$. Then $s_{i+1}=(s_i)_{d_{i+1}}$ is defined for $i<n$. Moreover $B_i(P)\subseteq A(s_i)$ for $0\le i\le n$, where the labels of $s_i$ are identified with physical directions through the canonicalizing symmetry.
\end{lemma}

\begin{proof}
By induction the retained suffix of $s_i$, transported back to the physical frame, is a final segment of the lift $(x_0,\dots,x_i)$. If $d\notin A(s_i)$, then $x_i+d$ is a retained vertex $x_j$, $j<i$. Its projection $v_j$ is visited, so $d\notin B_i(P)$. In particular $d_{i+1}\in B_i(P)\subseteq A(s_i)$, so the update is defined.
\end{proof}

Periodic identifications, and returns to vertices older than the retained suffix, can remove further directions from $B_i(P)$. The bounds below are designed so that removing directions only decreases them. They are therefore uniform in $L$.

\subsection{Transfer bounds}

Let $\varphi$ assign a weight $\varphi(s,d)\ge0$ to every labelled transition and let $\kappa>0$. For a vector $v$ indexed by states put
\begin{equation}\label{eq:transfer}
 (\cT v)(s)=\max_{B\subseteq A(s)}\kappa^{5-|B|}\sum_{d\in B}\varphi(s,d)\,v(s_d)\quad(s\ne s_0),\qquad
 (\cT v)(s_0)=\sum_{d\in A(s_0)}\varphi(s_0,d)\,v((s_0)_d).
\end{equation}
For $v\ge0$ and fixed $|B|$, the maximum is attained by the $|B|$ largest summands, counted with multiplicity. If $\kappa\le1$, it is attained at $B=A(s)$, and $\cT$ is linear.

\begin{lemma}[Transfer bound]\label{lem:transfer}
Suppose $v>0$ and $\cT v\le\Lambda v$ componentwise. For a path $P$ write $\Phi(P)=\prod_{i=0}^{n-1}\varphi(s_i,d_{i+1})$. Then for every even $L\ge8$ and $n\ge1$,
\[
 \sum_{|P|=n}\kappa^{c(P)}\,\Phi(P)\ \le\ \max(1,\kappa^5)\,\frac{v(s_0)}{\min_sv(s)}\,\Lambda^n .
\]
\end{lemma}

\begin{proof}
For a path $p$ of length $\ell$ with state $s(p)$, and $k\ge0$, let
\[
 F_k(p)=\sum_q\ \prod_{i=\ell}^{\ell+k-1}\kappa^{a_i(q)\one\{i\ge1\}}\varphi(s_i(q),d_{i+1}(q)),
\]
summed over all extensions $q$ of $p$ of length $\ell+k$. We show $F_k(p)\le\Lambda^kv(s(p))/\min v$. This is clear for $k=0$. Write $s=s(p)$ and $B=B_\ell(p)$, which satisfies $\kappa^{a_\ell}=\kappa^{5-|B|}$ for $\ell\ge1$ by \cref{lem:contacts}. Then
\[
\begin{aligned}
 F_{k+1}(p)&=\kappa^{(5-|B|)\one\{\ell\ge1\}}\sum_{d\in B}\varphi(s,d)F_k(p+d)\\
 &\le\frac{\Lambda^k}{\min v}\,\kappa^{(5-|B|)\one\{\ell\ge1\}}\sum_{d\in B}\varphi(s,d)v(s_d)
 \le\frac{\Lambda^k(\cT v)(s)}{\min v}.
\end{aligned}
\]
The first inequality uses the induction hypothesis and $s(p+d)=s_d$. The second uses $B\subseteq A(s)$ from \cref{lem:automaton} and the definition of $\cT$. Finally $(\cT v)(s)\le\Lambda v(s)$. Since $c(P)=\sum_{i=1}^na_i$, we get $\sum_{|P|=n}\kappa^{c(P)}\Phi(P)\le\max_Pw^{a_n(P)}F_n(\text{root})$, and $a_n\le5$.
\end{proof}

We use two instances of \cref{lem:transfer}. The \emph{contact operator} has $\varphi\equiv1$ and $\kappa=\kappa_1:=10000/9991=(1-\theta_1)^{-1}$. Since $e^{\theta_1}<\kappa_1$, it bounds $\sum_Pe^{\theta_1c(P)}$. The \emph{shape operator} has $\kappa=1$, $\varphi(s,d)=\overline\Omega_\tau$ if $s$ has at least two steps and $\tau$ is the shape of the triple completed by $d$, and $\varphi(s,d)=1$ otherwise. For it
\[
 \Phi(P)=\prod_{i=1}^{n-2}\overline\Omega_{\tau_i}\ \ge\ \Omega(P):=\exp\Bigl(-\sum_{i=1}^{n-2}w_{\tau_i}\Bigr),
\]
because $\overline\Omega_\tau\ge1-w_\tau+w_\tau^2/2\ge e^{-w_\tau}$ for $w_\tau\ge0$.

\begin{certificate}[Transfer bounds]\label{cert:transfer}
The automaton of \cref{def:automaton} has $12\,994\,574$ states, including the initial state, and $53\,221\,763$ labelled transitions.
\begin{enumerate}[label=\textup{(\alph*)}]
\item For the contact operator, with each power $\kappa_1^j$, $0\le j\le5$, rounded up to the grid $10^{-8}\Z$, there is a positive integer vector $v$ with $\cT v\le\Lambda_1v$ and
\[
 \Lambda_1=\frac{473271745}{10^8}.
\]
The largest row quotient is $7337944694907297/1550471750000000<\Lambda_1$.
\item For the shape operator, with the weights $\overline\Omega_\tau$ of \cref{tab:shapes}, there is a positive integer vector $v$ with $\cT v\le U_2v$ and
\[
 U_2=\frac{1182583}{250000}=4.730332 .
\]
\item For the \emph{unweighted operator}, with $\varphi\equiv1$ and $\kappa=1$, there is a positive integer vector $v$ with $\cT v\le U_0v$ and
\[
 U_0=\frac{473195}{100000}=4.73195 .
\]
The largest row quotient is $151907537380697/32102539197381<U_0$. The root counts of the automaton agree with the numbers of self-avoiding walks on $\Z^3$ of length $n\le17$~\cite{SBB11}.
\end{enumerate}
\end{certificate}

The vectors were obtained by iterating upward-rounded versions of $\cT$ from a constant vector, renormalizing by a common power of two, and adding one. How they were found plays no role. The certificate is the row-by-row inequality for the final vector. That inequality is checked with exact integer cross-multiplication.

\subsection{Combining shapes and chords}

The shape operator ignores contacts. The chord price is restored by interpolating between it and \cref{lem:rosenbluth}.

\begin{lemma}[H\"older combination]\label{lem:holder}
Let $t=\theta_2/\log(5/4)$. Then $0<t<1$, $(5/U_2)^4<5/4$, and for every even $L\ge8$ and $n\ge1$,
\[
 \sum_{|P|=n}\Omega(P)\,e^{\theta_2c(P)}\ \le\ C_2\bigl(U_2e^{\theta_2/4}\bigr)^n,
 \qquad C_2=\Bigl(\frac{v(s_0)}{\min v}\Bigr)^{1-t}\Bigl(\frac65\Bigl(\frac54\Bigr)^5\Bigr)^t,
\]
where $v$ is the vector of \cref{cert:transfer}\textup{(b)}.
\end{lemma}

\begin{proof}
The two numerical claims are elementary; see \cref{cert:endpoint-lower}. Write $e^{\theta_2c}=(5/4)^{tc}$ and $\Omega=\Omega^{1-t}\Omega^t$. H\"older's inequality with exponents $1/(1-t)$ and $1/t$ gives
\[
 \sum_P\Omega(P)e^{\theta_2c(P)}\le\Bigl(\sum_P\Omega(P)\Bigr)^{1-t}\Bigl(\sum_P\Omega(P)(5/4)^{c(P)}\Bigr)^t .
\]
By \cref{lem:transfer} for the shape operator, the first sum is at most $(v(s_0)/\min v)\,U_2^n$. Since $0<\Omega\le1$, \cref{lem:rosenbluth} bounds the second sum by $6(5/4)^55^{n-1}$. Hence the left-hand side is at most $C_2U_2^n(5/U_2)^{tn}$. Finally $(5/U_2)^t\le e^{\theta_2/4}$ is equivalent to $t\log(5/U_2)\le\frac t4\log\frac54$, i.e.\ to $(5/U_2)^4\le5/4$.
\end{proof}

\begin{remark}
The exponent $1/4$ in $e^{\theta_2/4}$ reflects how close $5$ is to $U_2$. The chord moment of \cref{lem:rosenbluth} grows like $5^n$ rather than $U_2^n$. The H\"older combination therefore costs a factor $(5/U_2)^t$ per step, and the inequality $(5/U_2)^4<5/4$ converts this into $e^{\theta_2/4}$. Its margin is small: $(5/U_2)^4=1.24828\ldots$.
\end{remark}
